\documentclass[aspectratio=169,10pt,english]{beamer}

\input{../../../_preambulo_beamer.tex}
\input{../../../_comandos_beamer.tex}

\selectlanguage{english}

\title{MA1001B - Statistical Modeling for Decision Making}
\subtitle{Lesson 01: Population, Samples, and Data Types}
\author{}
\institute{Tecnol\'ogico de Monterrey}
\date{}

\begin{document}

\begin{frame}[plain]
  \vspace{1.2cm}
  {\Large\bfseries MA1001B - Statistical Modeling for Decision Making\par}
  \vspace{0.7cm}
  {\large Lesson 01: Population, Samples, and Data Types\par}
  \vspace{0.7cm}
  {\color{TecRojo}\rule{\textwidth}{1.2pt}}
  \vspace{0.8cm}

  MA1001B Faculty\\[0.3cm]
  Term\\[0.3cm]
  Tecnol\'ogico de Monterrey
\end{frame}

\begin{frame}{The Statistical Question}
  \begin{alertblock}{Guiding question}
    How can a limited set of observations support a claim about a larger
    business population?
  \end{alertblock}

  \vspace{0.6em}
  By the end of this lesson, you can:
  \begin{itemize}
    \item define the target population and the observed sample;
    \item connect a population parameter with its sample statistic;
    \item classify variables from their analytical meaning;
    \item identify why sample size alone cannot remove selection bias.
  \end{itemize}

  \vfill
  \scriptsize All records used today are synthetic. No real company data are used.
\end{frame}

\begin{frame}{Population, Sample, Parameter, and Statistic}
  \small
  \begin{center}
    \begin{tabular}{p{0.19\textwidth}p{0.32\textwidth}p{0.32\textwidth}}
      \toprule
      \textbf{Role} & \textbf{Population} & \textbf{Sample} \\
      \midrule
      Observations & All units defined by the question & Units actually selected \\
      Size & $N$ & $n$ \\
      Numerical summary & Parameter, such as $\mu$ & Statistic, such as $\bar{x}$ \\
      Status & Fixed for the defined population & Changes across samples \\
      \bottomrule
    \end{tabular}
  \end{center}

  \vspace{0.6em}
  \begin{exampleblock}{Synthetic operations question}
    What is the mean processing time for all $12{,}000$ orders in the planning
    period? A random sample of 200 orders provides $\bar{x}$ as an estimate of
    the population mean $\mu$.
  \end{exampleblock}
\end{frame}

\begin{frame}[fragile]{Step 1: A Sample Estimates a Population Quantity}
  \begin{columns}[T]
    \begin{column}{0.50\textwidth}
      \small
      \begin{lstlisting}
population = build_population()
sample = draw_random_sample(population)

mu = population[variable].mean()
x_bar = sample[variable].mean()
error = abs(x_bar - mu)
      \end{lstlisting}

      \begin{block}{Verified output}
        $N=12{,}000$ and $n=200$\\
        $\mu=45.26$ minutes\\
        $\bar{x}=46.11$ minutes\\
        Absolute error $=0.85$ minutes
      \end{block}
    \end{column}
    \begin{column}{0.50\textwidth}
      \centering
      \includegraphics[width=\textwidth]{../figures/data_foundations_01_population_sample.png}
      \vspace{0.2em}
      \scriptsize Population and random-sample distributions from Step 1
    \end{column}
  \end{columns}
\end{frame}

\begin{frame}{Business Data Types}
  \small
  \begin{center}
    \begin{tabular}{p{0.23\textwidth}p{0.28\textwidth}p{0.34\textwidth}}
      \toprule
      \textbf{Type} & \textbf{Question} & \textbf{Lesson example} \\
      \midrule
      Categorical & Which group or label? & Channel, region \\
      Quantitative discrete & How many countable units? & Items per order \\
      \shortstack[l]{Quantitative\\continuous} & How much was measured? & Processing time, order value \\
      \bottomrule
    \end{tabular}
  \end{center}

  \vspace{0.7em}
  \begin{alertblock}{Analytical meaning controls the classification}
    \texttt{order\_id} contains digits, but arithmetic on an identifier has no
    business meaning. It acts as a label rather than a quantitative measurement.
  \end{alertblock}
\end{frame}

\begin{frame}[fragile]{Step 2: One Table Contains Several Data Types}
  \begin{columns}[T]
    \begin{column}{0.42\textwidth}
      \small
      \begin{lstlisting}
roles = {
  "channel": "Categorical",
  "items_per_order":
      "Quantitative discrete",
  "processing_time_minutes":
      "Quantitative continuous",
}
      \end{lstlisting}

      The same 200-row sample supports different summaries because each variable
      answers a different kind of question.
    \end{column}
    \begin{column}{0.58\textwidth}
      \centering
      \includegraphics[width=\textwidth]{../figures/data_foundations_02_data_types.png}
      \vspace{0.2em}
      \scriptsize Counts suit categories and discrete values; a histogram suits
      a measured continuous value.
    \end{column}
  \end{columns}
\end{frame}

\begin{frame}{Step 3: A Larger Biased Sample Has More Error Here}
  \begin{columns}[T]
    \begin{column}{0.43\textwidth}
      \small
      Population parameter: $\mu=45.26$ minutes

      \vspace{0.5em}
      \begin{tabular}{lrr}
        \toprule
        \textbf{Sample} & \textbf{$n$} & \textbf{Error} \\
        \midrule
        Random & 200 & 0.85 min \\
        Convenience & 1,000 & 6.24 min \\
        \bottomrule
      \end{tabular}

      \vspace{0.7em}
      The convenience sample contains only readily available Web orders. Its
      larger size does not restore the missing channels.

      \vspace{0.5em}
      \textbf{Limit:} This deterministic example demonstrates possible selection
      bias. It does not claim that every convenience sample has this exact error.
    \end{column}
    \begin{column}{0.57\textwidth}
      \centering
      \includegraphics[width=\textwidth]{../figures/data_foundations_03_sampling_bias.png}
      \vspace{0.2em}
      \scriptsize Sample size and estimation error from the same synthetic population
    \end{column}
  \end{columns}
\end{frame}

\begin{frame}{Concept Check}
  \begin{enumerate}
    \item An analyst studies all orders completed this month and inspects 300 of
      them. Identify the population and the sample.
    \item The 300 orders have a mean value of MXN 1,840. Is this value a
      parameter or a statistic? Explain why.
    \item Classify customer segment, number of items, and delivery time.
    \item Why can 1,000 conveniently available records provide a worse estimate
      than 200 randomly selected records?
  \end{enumerate}

  \vspace{0.6em}
  \begin{block}{Connection to Lesson 02}
    The session can continue directly into simple random, stratified, and cluster
    sampling designs.
  \end{block}

  \vfill
  \scriptsize Source: OpenStax, \textit{Introductory Business Statistics 2e},
  Chapter 1, Sections 1.1 and 1.2. CC BY-NC-SA.
\end{frame}

\appendix



\end{document}
